%%========================================================================
%% IDENTITAS TUGAS AKHIR
%% Isi semua bagian di bawah ini. Bagian yang masih contoh WAJIB diganti.
%%========================================================================

%--- Judul -----------------------------------------------------------------
% Ditulis singkat (sekitar 20 kata), kapital, membentuk piramida terbalik.
% Di LaTeX cukup tulis biasa; pemenggalan baris diatur otomatis.
\providecommand{\judulid}{Rancang Bangun Sistem Deteksi Kesegaran Ikan Berbasis Pengolahan Citra Digital dan Jaringan Saraf Tiruan}
\providecommand{\judulen}{Design of a Fish Freshness Detection System Based on Digital Image Processing and Artificial Neural Networks}

%--- Mahasiswa -------------------------------------------------------------
\providecommand{\penulis}{Nama Lengkap Mahasiswa}
\providecommand{\nim}{3210xxxxxx}

%--- Program studi ---------------------------------------------------------
\providecommand{\jenjang}{Sarjana Terapan}          % atau: Diploma Tiga
\providecommand{\gelar}{Sarjana Terapan (S.Tr.)}    % atau: Ahli Madya (A.Md)
\providecommand{\prodi}{Teknologi Rekayasa Perangkat Lunak}
\providecommand{\jurusan}{Teknik Informatika}
\providecommand{\institusi}{Politeknik Negeri Banyuwangi}
\providecommand{\kota}{Banyuwangi}
\providecommand{\tahun}{\the\year}

%--- Pembimbing dan penguji ------------------------------------------------
\providecommand{\pembimbingA}{Nama Pembimbing I, S.Kom., M.Kom.}
\providecommand{\NIDNpembimbingA}{00xxxxxxxx}
\providecommand{\pembimbingB}{Nama Pembimbing II, S.T., M.T.}
\providecommand{\NIDNpembimbingB}{00xxxxxxxx}

\providecommand{\pengujiA}{Nama Penguji I, S.Kom., M.Kom.}
\providecommand{\pengujiB}{Nama Penguji II, S.T., M.T.}

%--- Pejabat penanda tangan ------------------------------------------------
\providecommand{\ketuajurusan}{Nama Ketua Jurusan, S.T., M.T.}
\providecommand{\NIPketuajurusan}{19xxxxxxxxxxxxxxxx}
\providecommand{\kaprodi}{Nama Ketua Program Studi, S.Kom., M.Kom.}
\providecommand{\NIPkaprodi}{19xxxxxxxxxxxxxxxx}

%--- Pelaksanaan seminar / sidang ------------------------------------------
\providecommand{\hariujian}{Senin}
\providecommand{\tanggalUjian}{1 Januari \tahun}
\providecommand{\tempatujian}{Ruang Sidang Jurusan \jurusan}

%--- Warna sampul sesuai program studi (Pedoman Tabel 6.2) -----------------
% Teknik Informatika   : kuning, tulisan hitam   -> FFFF00
% Teknik Sipil         : biru tua, tulisan emas  -> 000080
% Teknik Mesin         : merah, tulisan emas     -> FF0000
% Manajemen Bisnis Pariwisata : oranye           -> FF7F00
% Agribisnis           : hijau tua               -> 009F3C
% Teknologi Pengolahan Hasil Ternak : coklat tua -> 964B00
% Teknik Manufaktur Kapal : biru muda            -> 0080FF
\providecommand{\warnasampul}{FFFF00}
